% Options for packages loaded elsewhere
\PassOptionsToPackage{unicode}{hyperref}
\PassOptionsToPackage{hyphens}{url}
\documentclass[
  nobib]{tufte-handout}
\usepackage{xcolor}
\usepackage{amsmath,amssymb}
\setcounter{secnumdepth}{-\maxdimen} % remove section numbering
\usepackage{iftex}
\ifPDFTeX
  \usepackage[T1]{fontenc}
  \usepackage[utf8]{inputenc}
  \usepackage{textcomp} % provide euro and other symbols
\else % if luatex or xetex
  \usepackage{unicode-math} % this also loads fontspec
  \defaultfontfeatures{Scale=MatchLowercase}
  \defaultfontfeatures[\rmfamily]{Ligatures=TeX,Scale=1}
\fi
\usepackage{lmodern}
\ifPDFTeX\else
  % xetex/luatex font selection
\fi
% Use upquote if available, for straight quotes in verbatim environments
\IfFileExists{upquote.sty}{\usepackage{upquote}}{}
\IfFileExists{microtype.sty}{% use microtype if available
  \usepackage[]{microtype}
  \UseMicrotypeSet[protrusion]{basicmath} % disable protrusion for tt fonts
}{}
\makeatletter
\@ifundefined{KOMAClassName}{% if non-KOMA class
  \IfFileExists{parskip.sty}{%
    \usepackage{parskip}
  }{% else
    \setlength{\parindent}{0pt}
    \setlength{\parskip}{6pt plus 2pt minus 1pt}}
}{% if KOMA class
  \KOMAoptions{parskip=half}}
\makeatother
\usepackage{color}
\usepackage{fancyvrb}
\newcommand{\VerbBar}{|}
\newcommand{\VERB}{\Verb[commandchars=\\\{\}]}
\DefineVerbatimEnvironment{Highlighting}{Verbatim}{commandchars=\\\{\}}
% Add ',fontsize=\small' for more characters per line
\newenvironment{Shaded}{}{}
\newcommand{\AlertTok}[1]{\textcolor[rgb]{1.00,0.00,0.00}{\textbf{#1}}}
\newcommand{\AnnotationTok}[1]{\textcolor[rgb]{0.38,0.63,0.69}{\textbf{\textit{#1}}}}
\newcommand{\AttributeTok}[1]{\textcolor[rgb]{0.49,0.56,0.16}{#1}}
\newcommand{\BaseNTok}[1]{\textcolor[rgb]{0.25,0.63,0.44}{#1}}
\newcommand{\BuiltInTok}[1]{\textcolor[rgb]{0.00,0.50,0.00}{#1}}
\newcommand{\CharTok}[1]{\textcolor[rgb]{0.25,0.44,0.63}{#1}}
\newcommand{\CommentTok}[1]{\textcolor[rgb]{0.38,0.63,0.69}{\textit{#1}}}
\newcommand{\CommentVarTok}[1]{\textcolor[rgb]{0.38,0.63,0.69}{\textbf{\textit{#1}}}}
\newcommand{\ConstantTok}[1]{\textcolor[rgb]{0.53,0.00,0.00}{#1}}
\newcommand{\ControlFlowTok}[1]{\textcolor[rgb]{0.00,0.44,0.13}{\textbf{#1}}}
\newcommand{\DataTypeTok}[1]{\textcolor[rgb]{0.56,0.13,0.00}{#1}}
\newcommand{\DecValTok}[1]{\textcolor[rgb]{0.25,0.63,0.44}{#1}}
\newcommand{\DocumentationTok}[1]{\textcolor[rgb]{0.73,0.13,0.13}{\textit{#1}}}
\newcommand{\ErrorTok}[1]{\textcolor[rgb]{1.00,0.00,0.00}{\textbf{#1}}}
\newcommand{\ExtensionTok}[1]{#1}
\newcommand{\FloatTok}[1]{\textcolor[rgb]{0.25,0.63,0.44}{#1}}
\newcommand{\FunctionTok}[1]{\textcolor[rgb]{0.02,0.16,0.49}{#1}}
\newcommand{\ImportTok}[1]{\textcolor[rgb]{0.00,0.50,0.00}{\textbf{#1}}}
\newcommand{\InformationTok}[1]{\textcolor[rgb]{0.38,0.63,0.69}{\textbf{\textit{#1}}}}
\newcommand{\KeywordTok}[1]{\textcolor[rgb]{0.00,0.44,0.13}{\textbf{#1}}}
\newcommand{\NormalTok}[1]{#1}
\newcommand{\OperatorTok}[1]{\textcolor[rgb]{0.40,0.40,0.40}{#1}}
\newcommand{\OtherTok}[1]{\textcolor[rgb]{0.00,0.44,0.13}{#1}}
\newcommand{\PreprocessorTok}[1]{\textcolor[rgb]{0.74,0.48,0.00}{#1}}
\newcommand{\RegionMarkerTok}[1]{#1}
\newcommand{\SpecialCharTok}[1]{\textcolor[rgb]{0.25,0.44,0.63}{#1}}
\newcommand{\SpecialStringTok}[1]{\textcolor[rgb]{0.73,0.40,0.53}{#1}}
\newcommand{\StringTok}[1]{\textcolor[rgb]{0.25,0.44,0.63}{#1}}
\newcommand{\VariableTok}[1]{\textcolor[rgb]{0.10,0.09,0.49}{#1}}
\newcommand{\VerbatimStringTok}[1]{\textcolor[rgb]{0.25,0.44,0.63}{#1}}
\newcommand{\WarningTok}[1]{\textcolor[rgb]{0.38,0.63,0.69}{\textbf{\textit{#1}}}}
\usepackage{longtable,booktabs,array}
\newcounter{none} % for unnumbered tables
\usepackage{calc} % for calculating minipage widths
% Correct order of tables after \paragraph or \subparagraph
\usepackage{etoolbox}
\makeatletter
\patchcmd\longtable{\par}{\if@noskipsec\mbox{}\fi\par}{}{}
\makeatother
% Allow footnotes in longtable head/foot
\IfFileExists{footnotehyper.sty}{\usepackage{footnotehyper}}{\usepackage{footnote}}
\makesavenoteenv{longtable}
\setlength{\emergencystretch}{3em} % prevent overfull lines
\providecommand{\tightlist}{%
  \setlength{\itemsep}{0pt}\setlength{\parskip}{0pt}}
% =====================================================================
% axona-doc-preamble.tex — shared LaTeX preamble for the programmer-guide
% quartet (Quick-Start, Programmer Guide, API Reference, Services Guide).
%
% These docs are markdown-native; the PDF is built by converting the .md to
% a tufte-handout .tex with pandoc (-H this file) and compiling with tectonic.
% See render.sh. Modeled on explainer/Axona-Explainer.tex (the flat-essay
% original); the tweaks below make tufte-handout work for code/table-heavy
% reference material.
% =====================================================================

\definecolor{rust}{RGB}{185,78,54}
\definecolor{rustsoft}{RGB}{212,168,150}
\definecolor{inkmuted}{RGB}{102,102,102}
\usepackage{microtype}

% Explainer-matched page furniture: tufte-handout's letterspaced running
% heads fail on certain page breaks (the SOUL \MakeTextLowercase issue the
% explainer works around), so use the same fancyhdr style — no running
% head, a centered page number in the footer.
\usepackage{fancyhdr}
\fancypagestyle{ndhtplain}{%
  \fancyhf{}%
  \renewcommand{\headrulewidth}{0pt}%
  \fancyfoot[C]{\thepage}%
}

% Explainer-matched heading-orphan protection: reserve vertical room before
% each (sub)section heading so a heading never sits alone at a page bottom.
\usepackage{etoolbox}
\usepackage{needspace}
\pretocmd{\section}{\needspace{6\baselineskip}}{}{}
\pretocmd{\subsection}{\needspace{4\baselineskip}}{}{}

% XeTeX (tectonic) unicode monospace: renders … → and the box-drawing ASCII
% diagrams that the default CM/LM mono cannot. Menlo ships with macOS.
\setmonofont{Menlo}[Scale=0.78]

% tufte-handout ships \subsubsection as an error stub (it wants flat essays).
% A run-in fallback lets nested reference headings (§4.2, §12.3, …) compile.
\renewcommand{\subsubsection}[1]{\medskip\par\noindent{\normalsize\bfseries #1}\par\nobreak\smallskip}

% These docs use no margin notes, so reclaim tufte's wide right margin: widen
% the text block so code blocks and API tables don't overflow it.
\usepackage{geometry}
\geometry{left=1in, textwidth=6.1in, marginparsep=0.2in, marginparwidth=0.9in,
          top=1in, headsep=0.4in}

% Wrap long code lines (cheat-sheet comments run to ~130 chars) with a ↪ marker
% instead of overflowing. Applies to pandoc's fancyvrb Highlighting/Verbatim.
\usepackage{fvextra}
\fvset{breaklines=true, breakanywhere=true}

% A few technical symbols appear in prose / inline code that Latin Modern and
% Menlo don't cover; map them to math equivalents so they render (rather than
% dropping to tofu). Only maps chars actually used — leaves → alone (it renders).
\usepackage{newunicodechar}
\newunicodechar{≈}{\ensuremath{\approx}}
\newunicodechar{≠}{\ensuremath{\neq}}
\newunicodechar{≥}{\ensuremath{\geq}}
\newunicodechar{≤}{\ensuremath{\leq}}
\newunicodechar{×}{\ensuremath{\times}}
\newunicodechar{·}{\ensuremath{\cdot}}
\newunicodechar{‖}{\ensuremath{\|}}
\usepackage{bookmark}
\IfFileExists{xurl.sty}{\usepackage{xurl}}{} % add URL line breaks if available
\urlstyle{same}
\hypersetup{
  hidelinks,
  pdfcreator={LaTeX via pandoc}}

\author{}
\date{}

\begin{document}

% Generated by render.sh — the explainer-style title page.
\thispagestyle{empty}
\null\vfill
\begin{center}
{\fontsize{48}{52}\selectfont \textsc{Axona}\par}
\vspace{0.6em}
{\Large\itshape Building Applications That Have No Server\par}
\vspace{4em}
{\large David A.~Smith\par}
{\itshape\small Axona.net\par}
\vspace{2em}
{\small The Axona Programmer Guide \quad v4.27.1 \quad 2026-07-17\par}
\end{center}
\vfill
\null
\newpage
\pagestyle{ndhtplain}
\thispagestyle{ndhtplain}
\setcounter{page}{1}

\section{The Axona Programmer Guide}\label{the-axona-programmer-guide}

\textbf{Build apps that have no server.} \emph{(kernel 4.27.1 ·
testnet)}

You're about to build something unusual: an application with no backend.
No database to stand up, no message broker to rent, no API keys, nothing
to deploy before your first user shows up. Your users' browsers and
devices \emph{are} the infrastructure --- they find each other, carry
each other's messages, and keep the whole thing alive.

This guide is for application programmers. You do not need to know how
distributed hash tables work, and this guide will not make you learn.
(If the machinery genuinely interests you, it's all in
\hyperref[9-under-the-hood--optional]{Under the hood} at the end ---
strictly optional.)

\textbf{Companions:}

\begin{itemize}
\tightlist
\item
  \href{Quick-Start-v4.27.1.md}{Quick Start} --- a working roundtrip in
  5 minutes. Do it first.
\item
  \href{Axona-API-Reference-v4.27.1.md}{API Reference} --- every
  signature, exactly.
\item
  \href{Axona-AI-Grounding-v4.27.1.md}{AI Grounding} --- building with
  an AI assistant? Hand it this file.
\item
  \href{Axona-Services-Guide-v4.27.1.md}{Services Guide} --- running
  bridges and relays (you can skip this for a long time).
\end{itemize}

\begin{quote}
\textbf{Building with an AI?} Most Axona apps are. Paste
\url{Axona-AI-Grounding-v4.27.1.md} into your assistant's context and it
will know every rule in this guide. Keep this guide for \emph{you} ---
it explains the why.
\end{quote}

\subsection{Table of contents}\label{table-of-contents}

\begin{enumerate}
\def\labelenumi{\arabic{enumi}.}
\tightlist
\item
  \hyperref[1-the-five-ideas]{The five ideas}
\item
  \hyperref[2-first-contact]{First contact}
\item
  \hyperref[3-idea-by-idea-gently]{Idea by idea, gently}
\item
  \hyperref[4-recipes]{Recipes}
\item
  \hyperref[5-things-that-will-bite-you]{Things that will bite you}
\item
  \hyperref[6-errors-when-they-happen]{Errors, when they happen}
\item
  \hyperref[7-limits-worth-knowing]{Limits worth knowing}
\item
  \hyperref[8-going-live]{Going live}
\item
  \hyperref[9-under-the-hood--optional]{Under the hood}
  \emph{(optional)}
\item
  \hyperref[10-appendix]{Appendix: what's coming, and migrating from
  older lines}
\end{enumerate}

\begin{center}\rule{0.5\linewidth}{0.5pt}\end{center}

\subsection{1. The five ideas}\label{the-five-ideas}

Everything you'll ever do with Axona is a combination of five ideas.
Here they are in full; the rest of this guide is elaboration.

\begin{enumerate}
\def\labelenumi{\arabic{enumi}.}
\item
  \textbf{A peer.} Your app running, connected to the network. It has a
  \emph{connection identity} minted fresh each run --- think of it as
  the seat your app happens to be sitting in today.
\item
  \textbf{An author.} A signing key that says \emph{who wrote this}. It
  has no location, works from any device, and is the only thing worth
  saving to disk. If your app has ``users,'' an author key is a user.
\item
  \textbf{A topic.} A place messages go. Not a string --- a small
  object:
  \texttt{\{\ region:\ \textquotesingle{}useast\textquotesingle{},\ name:\ \textquotesingle{}lobby\textquotesingle{}\ \}}.
  Anyone who writes the same fields gets the same place. No
  registration, no server assigning channels: shared fields \emph{are}
  the rendezvous.
\item
  \textbf{Publish.} Send a message to a topic, signed by an author:
  \texttt{peer.pub(topic,\ message,\ \{\ signWith:\ author\ \})}.
\item
  \textbf{Subscribe.} Receive a topic's messages, live and (if you ask)
  recent history:
  \texttt{peer.sub(topic,\ handler,\ \{\ since:\ \textquotesingle{}all\textquotesingle{}\ \})}.
\end{enumerate}

That's the protocol. Chat, feeds, presence, file sharing, multiplayer
state, agent-to-agent traffic --- they're all arrangements of these
five.

\subsection{2. First contact}\label{first-contact}

If you haven't run the \href{Quick-Start-v4.27.1.md}{Quick Start}, do it
now --- five minutes, a real message through the real network. The heart
of it:

\begin{Shaded}
\begin{Highlighting}[]
\CommentTok{// after connect() (Quick Start step 4; one call — see §3.1)}

\KeywordTok{const}\NormalTok{ TOPIC }\OperatorTok{=}\NormalTok{ \{ }\DataTypeTok{region}\OperatorTok{:} \StringTok{\textquotesingle{}useast\textquotesingle{}}\OperatorTok{,} \DataTypeTok{name}\OperatorTok{:} \StringTok{\textquotesingle{}quick{-}start{-}demo\textquotesingle{}}\NormalTok{ \}}\OperatorTok{;}

\ControlFlowTok{await}\NormalTok{ peer}\OperatorTok{.}\FunctionTok{sub}\NormalTok{(TOPIC}\OperatorTok{,}\NormalTok{ (env) }\KeywordTok{=\textgreater{}}\NormalTok{ \{}
  \BuiltInTok{console}\OperatorTok{.}\FunctionTok{log}\NormalTok{(}\StringTok{\textquotesingle{}[recv]\textquotesingle{}}\OperatorTok{,}\NormalTok{ env}\OperatorTok{.}\AttributeTok{message}\NormalTok{)}\OperatorTok{;}
\NormalTok{\}}\OperatorTok{,}\NormalTok{ \{ }\DataTypeTok{since}\OperatorTok{:} \StringTok{\textquotesingle{}all\textquotesingle{}}\NormalTok{ \})}\OperatorTok{;}

\ControlFlowTok{await}\NormalTok{ peer}\OperatorTok{.}\FunctionTok{pub}\NormalTok{(TOPIC}\OperatorTok{,} \StringTok{\textquotesingle{}hello, everyone\textquotesingle{}}\OperatorTok{,}\NormalTok{ \{ }\DataTypeTok{signWith}\OperatorTok{:}\NormalTok{ author \})}\OperatorTok{;}
\end{Highlighting}
\end{Shaded}

You publish; the network routes it; everyone subscribed --- including
you --- receives it. That echo of your own message is not an accident,
and you'll learn to love it (§5, ``There is no ack'').

\subsection{3. Idea by idea, gently}\label{idea-by-idea-gently}

\subsubsection{3.1 The peer --- your seat at the
table}\label{the-peer-your-seat-at-the-table}

A peer is created in one call, pointed at a \emph{bridge} --- the
meeting point where peers first find each other (after that, they talk
directly):

\begin{Shaded}
\begin{Highlighting}[]
\ImportTok{import}\NormalTok{ \{ connect \} }\ImportTok{from} \StringTok{\textquotesingle{}@axona/protocol/connect.js\textquotesingle{}}\OperatorTok{;}

\KeywordTok{const}\NormalTok{ \{ peer}\OperatorTok{,}\NormalTok{ author}\OperatorTok{,}\NormalTok{ status}\OperatorTok{,}\NormalTok{ disconnect \} }\OperatorTok{=} \ControlFlowTok{await} \FunctionTok{connect}\NormalTok{(\{}
  \DataTypeTok{bridge}\OperatorTok{:}   \StringTok{\textquotesingle{}wss://testnet.axona.net\textquotesingle{}}\OperatorTok{,}
  \DataTypeTok{location}\OperatorTok{:}\NormalTok{ \{ }\DataTypeTok{lat}\OperatorTok{:} \FloatTok{38.0}\OperatorTok{,} \DataTypeTok{lng}\OperatorTok{:} \OperatorTok{{-}}\FloatTok{77.0}\NormalTok{ \}}\OperatorTok{,}   \CommentTok{// the user\textquotesingle{}s real location}
  \DataTypeTok{author}\OperatorTok{:}   \StringTok{\textquotesingle{}myapp:author\textquotesingle{}}\OperatorTok{,}             \CommentTok{// persist the author key under this name}
\NormalTok{\})}\OperatorTok{;}
\CommentTok{// status: \{ ready: true, peers: 13, ms: 452, reason: \textquotesingle{}minPeers\textquotesingle{} \}}
\end{Highlighting}
\end{Shaded}

That's it --- identities minted, transport connected, mesh warmed. When
your app closes, \texttt{await\ disconnect()} leaves gracefully.

Under the sugar --- the primitives, if you need custom wiring

\texttt{connect()} composes public kernel pieces you can always use
directly (a custom transport, several peers on one page, a persistence
adapter):

\begin{Shaded}
\begin{Highlighting}[]
\ImportTok{import}\NormalTok{ \{}
\NormalTok{  AxonaPeer}\OperatorTok{,}\NormalTok{ AxonaDomain}\OperatorTok{,}\NormalTok{ NeuronNode}\OperatorTok{,}
\NormalTok{  createNodeIdentity}\OperatorTok{,}\NormalTok{ createAuthorIdentity}\OperatorTok{,}
\NormalTok{\} }\ImportTok{from} \StringTok{\textquotesingle{}@axona/protocol\textquotesingle{}}\OperatorTok{;}
\ImportTok{import}\NormalTok{ \{ webTransport \} }\ImportTok{from} \StringTok{\textquotesingle{}@axona/protocol/transport/web/index.js\textquotesingle{}}\OperatorTok{;}

\KeywordTok{const}\NormalTok{ nodeIdentity }\OperatorTok{=} \ControlFlowTok{await} \FunctionTok{createNodeIdentity}\NormalTok{(\{ }\DataTypeTok{lat}\OperatorTok{:} \FloatTok{38.0}\OperatorTok{,} \DataTypeTok{lng}\OperatorTok{:} \OperatorTok{{-}}\FloatTok{77.0}\NormalTok{ \})}\OperatorTok{;}
\KeywordTok{const}\NormalTok{ author       }\OperatorTok{=} \ControlFlowTok{await} \FunctionTok{createAuthorIdentity}\NormalTok{(\{ }\DataTypeTok{persistAs}\OperatorTok{:} \StringTok{\textquotesingle{}myapp:author\textquotesingle{}}\NormalTok{ \})}\OperatorTok{;}
\KeywordTok{const}\NormalTok{ transport    }\OperatorTok{=} \FunctionTok{webTransport}\NormalTok{(\{ }\DataTypeTok{bridgeUrl}\OperatorTok{:} \StringTok{\textquotesingle{}wss://testnet.axona.net\textquotesingle{}}\OperatorTok{,}
                                    \DataTypeTok{identity}\OperatorTok{:}\NormalTok{ nodeIdentity \})}\OperatorTok{;}
\KeywordTok{const}\NormalTok{ node }\OperatorTok{=} \KeywordTok{new} \FunctionTok{NeuronNode}\NormalTok{(\{ }\DataTypeTok{id}\OperatorTok{:}\NormalTok{ nodeIdentity}\OperatorTok{.}\AttributeTok{id}\OperatorTok{,} \DataTypeTok{lat}\OperatorTok{:} \FloatTok{38.0}\OperatorTok{,} \DataTypeTok{lng}\OperatorTok{:} \OperatorTok{{-}}\FloatTok{77.0}\NormalTok{ \})}\OperatorTok{;}
\NormalTok{node}\OperatorTok{.}\AttributeTok{transport} \OperatorTok{=}\NormalTok{ transport}\OperatorTok{;}
\KeywordTok{const}\NormalTok{ peer }\OperatorTok{=} \KeywordTok{new} \FunctionTok{AxonaPeer}\NormalTok{(\{ }\DataTypeTok{domain}\OperatorTok{:} \KeywordTok{new} \FunctionTok{AxonaDomain}\NormalTok{(\{ }\DataTypeTok{k}\OperatorTok{:} \DecValTok{20}\NormalTok{ \})}\OperatorTok{,}
\NormalTok{                             node}\OperatorTok{,}\NormalTok{ nodeIdentity}\OperatorTok{,}\NormalTok{ transport \})}\OperatorTok{;}
\ControlFlowTok{await}\NormalTok{ transport}\OperatorTok{.}\FunctionTok{start}\NormalTok{(nodeIdentity}\OperatorTok{.}\AttributeTok{id}\NormalTok{)}\OperatorTok{;}
\ControlFlowTok{await}\NormalTok{ peer}\OperatorTok{.}\FunctionTok{start}\NormalTok{()}\OperatorTok{;}
\ControlFlowTok{await}\NormalTok{ peer}\OperatorTok{.}\FunctionTok{ready}\NormalTok{()}\OperatorTok{;}
\end{Highlighting}
\end{Shaded}

Why the location? Axona is geographic: the first byte of every address
is a region of the Earth, and traffic for a region's topics stays among
that region's peers. Your user in Virginia is a \texttt{useast} peer;
their local topics never detour through Frankfurt.

The connection identity is deliberately disposable --- fresh each run,
like a dynamic IP. Don't save it, don't display it, don't build on it.
The thing with a lifespan is:

\subsubsection{3.2 The author --- the only key you
keep}\label{the-author-the-only-key-you-keep}

\begin{Shaded}
\begin{Highlighting}[]
\KeywordTok{const}\NormalTok{ author }\OperatorTok{=} \ControlFlowTok{await} \FunctionTok{createAuthorIdentity}\NormalTok{(\{ }\DataTypeTok{persistAs}\OperatorTok{:} \StringTok{\textquotesingle{}myapp:author\textquotesingle{}}\NormalTok{ \})}\OperatorTok{;}
\end{Highlighting}
\end{Shaded}

An author is an Ed25519 keypair. Its public half ---
\texttt{author.authorId} --- is what other people recognize: it appears
(as \texttt{signerPubkey}) on every message it signs, from any device,
from any location, forever. Persist it (\texttt{persistAs} uses
localStorage in the browser) and your user has a stable identity across
sessions. Skip \texttt{persistAs} and you have a throwaway persona ---
also sometimes exactly what you want. One app can hold several authors
at once (work self, home self, bot self).

There's no account creation, no email verification, no password reset
--- because there's no one to reset it \emph{with}. The key is the
identity. That's liberating and unforgiving in equal measure: lose the
key, lose the identity. Persist it.

\subsubsection{3.3 The topic --- an address you compute, not a channel
you
create}\label{the-topic-an-address-you-compute-not-a-channel-you-create}

\begin{Shaded}
\begin{Highlighting}[]
\KeywordTok{const}\NormalTok{ lobby }\OperatorTok{=}\NormalTok{ \{ }\DataTypeTok{region}\OperatorTok{:} \StringTok{\textquotesingle{}useast\textquotesingle{}}\OperatorTok{,} \DataTypeTok{name}\OperatorTok{:} \StringTok{\textquotesingle{}lobby\textquotesingle{}}\NormalTok{ \}}\OperatorTok{;}                        \CommentTok{// anyone writes}
\KeywordTok{const}\NormalTok{ feed  }\OperatorTok{=}\NormalTok{ \{ }\DataTypeTok{region}\OperatorTok{:} \StringTok{\textquotesingle{}useast\textquotesingle{}}\OperatorTok{,} \DataTypeTok{owner}\OperatorTok{:}\NormalTok{ me}\OperatorTok{.}\AttributeTok{authorId}\OperatorTok{,} \DataTypeTok{name}\OperatorTok{:} \StringTok{\textquotesingle{}posts\textquotesingle{}}\NormalTok{ \}}\OperatorTok{;}   \CommentTok{// only I write}
\KeywordTok{const}\NormalTok{ inbox }\OperatorTok{=}\NormalTok{ \{ }\DataTypeTok{region}\OperatorTok{:} \StringTok{\textquotesingle{}useast\textquotesingle{}}\OperatorTok{,} \DataTypeTok{owner}\OperatorTok{:}\NormalTok{ me}\OperatorTok{.}\AttributeTok{authorId}\OperatorTok{,} \DataTypeTok{name}\OperatorTok{:} \StringTok{\textquotesingle{}inbox\textquotesingle{}}\OperatorTok{,} \DataTypeTok{write}\OperatorTok{:} \StringTok{\textquotesingle{}open\textquotesingle{}}\NormalTok{ \}}\OperatorTok{;} \CommentTok{// anyone writes TO me}
\end{Highlighting}
\end{Shaded}

Those three shapes cover a startling share of application design:

\begin{itemize}
\tightlist
\item
  \textbf{The open room} (\texttt{\{\ region,\ name\ \}}) --- a chat, a
  game lobby, a public firehose. Anyone who knows the fields can read
  and write.
\item
  \textbf{The owned feed} (\texttt{owner} present; write defaults to
  \texttt{\textquotesingle{}owner\textquotesingle{}}) --- a blog, a
  profile, a bot's announcements. Everyone can read; only the owner's
  key can write, and the \emph{network itself} enforces that --- a
  forged post is refused by the peers storing the topic, not just by
  polite clients.
\item
  \textbf{The inbox} (\texttt{owner} +
  \texttt{write:\ \textquotesingle{}open\textquotesingle{}}) ---
  messages \emph{to} someone. Anyone writes; the owner's app reads.
\end{itemize}

Two rules keep topic bugs out of your life. First, \textbf{publisher and
subscriber must use identical fields} --- the topic ID is a hash of
them, so
\texttt{\{\ region:\ \textquotesingle{}useast\textquotesingle{},\ name:\ \textquotesingle{}lobby\textquotesingle{}\ \}}
and \texttt{\{\ name:\ \textquotesingle{}lobby\textquotesingle{}\ \}}
are different places (the second defaulted its region to wherever the
publisher happened to be). Centralize your descriptors in one function
and call it from both sides. Second, \textbf{prefer naming the region
explicitly} for anything two differently-located users must both find.

Need a shareable read handle? \texttt{await\ deriveTopicId(topic)}
returns the 66-hex ID; \texttt{sub} and \texttt{pull} accept it
directly.

\subsubsection{3.4 Publish --- signed, fire-and-forget, and that's a
feature}\label{publish-signed-fire-and-forget-and-thats-a-feature}

\begin{Shaded}
\begin{Highlighting}[]
\KeywordTok{const}\NormalTok{ msgId }\OperatorTok{=} \ControlFlowTok{await}\NormalTok{ peer}\OperatorTok{.}\FunctionTok{pub}\NormalTok{(topic}\OperatorTok{,}\NormalTok{ \{ }\DataTypeTok{text}\OperatorTok{:} \StringTok{\textquotesingle{}shipped it\textquotesingle{}}\NormalTok{ \}}\OperatorTok{,}\NormalTok{ \{ }\DataTypeTok{signWith}\OperatorTok{:}\NormalTok{ author \})}\OperatorTok{;}
\end{Highlighting}
\end{Shaded}

The message is any JSON-serializable value. The signature travels with
it; every receiver verifies it before your handler ever sees it. The
returned \texttt{msgId} is a \emph{content address} --- a hash of author
+ payload --- which buys you two elegant behaviors free of charge:
publishing the identical thing twice doesn't duplicate (it refreshes),
and any message can later be fetched or retracted by its id.

To publish anonymously, say so out loud:
\texttt{\{\ signWith:\ ANONYMOUS\ \}}. Forgetting \texttt{signWith}
entirely is an error --- Axona never guesses who's speaking.

And no, there is no delivery callback. See §5 for why that's a feature
and what to do instead (spoiler: you're already subscribed --- watch
your own message arrive).

\subsubsection{3.5 Subscribe --- live tail, with as much history as you
ask for}\label{subscribe-live-tail-with-as-much-history-as-you-ask-for}

\begin{Shaded}
\begin{Highlighting}[]
\KeywordTok{const}\NormalTok{ sub }\OperatorTok{=} \ControlFlowTok{await}\NormalTok{ peer}\OperatorTok{.}\FunctionTok{sub}\NormalTok{(topic}\OperatorTok{,}\NormalTok{ (env) }\KeywordTok{=\textgreater{}}\NormalTok{ \{}
  \ControlFlowTok{if}\NormalTok{ (env}\OperatorTok{.}\AttributeTok{deleted}\NormalTok{) }\ControlFlowTok{return} \FunctionTok{removeFromUI}\NormalTok{(env}\OperatorTok{.}\AttributeTok{msgId}\NormalTok{)}\OperatorTok{;}      \CommentTok{// a retraction marker}
  \FunctionTok{render}\NormalTok{(env}\OperatorTok{.}\AttributeTok{message}\OperatorTok{,}\NormalTok{ env}\OperatorTok{.}\AttributeTok{signerPubkey}\NormalTok{)}\OperatorTok{;}                \CommentTok{// verified author (or undefined)}
\NormalTok{\}}\OperatorTok{,}\NormalTok{ \{ }\DataTypeTok{since}\OperatorTok{:} \StringTok{\textquotesingle{}all\textquotesingle{}}\NormalTok{ \})}\OperatorTok{;}
\end{Highlighting}
\end{Shaded}

The \texttt{since} option is your replay dial:

{\def\LTcaptype{none} % do not increment counter
\begin{longtable}[]{@{}ll@{}}
\toprule\noalign{}
\texttt{since} & You receive \\
\midrule\noalign{}
\endhead
\bottomrule\noalign{}
\endlastfoot
\emph{(omitted)} & live messages only, from now \\
\texttt{\textquotesingle{}latest\textquotesingle{}} & the newest cached
message, then live \\
\texttt{\textquotesingle{}all\textquotesingle{}} & recent history (the
cached queue, \textasciitilde24 h), then live \\
\texttt{1719954000000} & everything newer than that timestamp, then
live \\
\end{longtable}
}

Deliveries are exactly-once per message per subscriber --- the kernel
dedups by \texttt{msgId}, so you never write defensive ``have I seen
this?'' code. \texttt{sub.stop()} ends one subscription;
\texttt{peer.unsub(topic)} ends all of yours for that topic.

\subsection{4. Recipes}\label{recipes}

Real shapes, ready to lift. Each assumes the assembled \texttt{peer} and
\texttt{author} from §3.1.

\subsubsection{4.1 A chat room}\label{a-chat-room}

\begin{Shaded}
\begin{Highlighting}[]
\KeywordTok{const}\NormalTok{ room }\OperatorTok{=}\NormalTok{ (name) }\KeywordTok{=\textgreater{}}\NormalTok{ (\{ }\DataTypeTok{region}\OperatorTok{:} \StringTok{\textquotesingle{}useast\textquotesingle{}}\OperatorTok{,} \DataTypeTok{name}\OperatorTok{:} \VerbatimStringTok{\textasciigrave{}chat:}\SpecialCharTok{$\{}\NormalTok{name}\SpecialCharTok{\}}\VerbatimStringTok{\textasciigrave{}}\NormalTok{ \})}\OperatorTok{;}

\KeywordTok{async} \KeywordTok{function} \FunctionTok{join}\NormalTok{(name}\OperatorTok{,}\NormalTok{ onMessage) \{}
  \ControlFlowTok{return}\NormalTok{ peer}\OperatorTok{.}\FunctionTok{sub}\NormalTok{(}\FunctionTok{room}\NormalTok{(name)}\OperatorTok{,}\NormalTok{ (env) }\KeywordTok{=\textgreater{}}\NormalTok{ \{}
    \ControlFlowTok{if}\NormalTok{ (env}\OperatorTok{.}\AttributeTok{deleted}\NormalTok{) }\ControlFlowTok{return}\OperatorTok{;}
    \FunctionTok{onMessage}\NormalTok{(\{ }\DataTypeTok{text}\OperatorTok{:}\NormalTok{ env}\OperatorTok{.}\AttributeTok{message}\OperatorTok{.}\AttributeTok{text}\OperatorTok{,}
                \DataTypeTok{from}\OperatorTok{:}\NormalTok{ (env}\OperatorTok{.}\AttributeTok{signerPubkey} \OperatorTok{??} \StringTok{\textquotesingle{}anon\textquotesingle{}}\NormalTok{)}\OperatorTok{.}\FunctionTok{slice}\NormalTok{(}\DecValTok{0}\OperatorTok{,} \DecValTok{8}\NormalTok{)}\OperatorTok{,}
                \DataTypeTok{at}\OperatorTok{:}\NormalTok{   env}\OperatorTok{.}\AttributeTok{ts}\NormalTok{ \})}\OperatorTok{;}
\NormalTok{  \}}\OperatorTok{,}\NormalTok{ \{ }\DataTypeTok{since}\OperatorTok{:} \StringTok{\textquotesingle{}all\textquotesingle{}}\NormalTok{ \})}\OperatorTok{;}          \CommentTok{// newcomers see recent history}
\NormalTok{\}}

\KeywordTok{async} \KeywordTok{function} \FunctionTok{say}\NormalTok{(name}\OperatorTok{,}\NormalTok{ text) \{}
  \ControlFlowTok{return}\NormalTok{ peer}\OperatorTok{.}\FunctionTok{pub}\NormalTok{(}\FunctionTok{room}\NormalTok{(name)}\OperatorTok{,}\NormalTok{ \{ text \}}\OperatorTok{,}\NormalTok{ \{ }\DataTypeTok{signWith}\OperatorTok{:}\NormalTok{ author \})}\OperatorTok{;}
\NormalTok{\}}
\end{Highlighting}
\end{Shaded}

That's a chat app. History for late joiners, live tail for everyone,
authorship on every line --- and notice what's missing: no room creation
step. The first person to publish \emph{is} the room springing into
existence.

\subsubsection{4.2 A personal feed (and reading someone
else's)}\label{a-personal-feed-and-reading-someone-elses}

\begin{Shaded}
\begin{Highlighting}[]
\KeywordTok{const}\NormalTok{ myFeed }\OperatorTok{=}\NormalTok{ \{ }\DataTypeTok{region}\OperatorTok{:} \StringTok{\textquotesingle{}useast\textquotesingle{}}\OperatorTok{,} \DataTypeTok{owner}\OperatorTok{:}\NormalTok{ author}\OperatorTok{.}\AttributeTok{authorId}\OperatorTok{,} \DataTypeTok{name}\OperatorTok{:} \StringTok{\textquotesingle{}posts\textquotesingle{}}\NormalTok{ \}}\OperatorTok{;}

\ControlFlowTok{await}\NormalTok{ peer}\OperatorTok{.}\FunctionTok{pub}\NormalTok{(myFeed}\OperatorTok{,}\NormalTok{ \{ }\DataTypeTok{title}\OperatorTok{:} \StringTok{\textquotesingle{}First post\textquotesingle{}}\OperatorTok{,} \DataTypeTok{body}\OperatorTok{:} \StringTok{\textquotesingle{}…\textquotesingle{}}\NormalTok{ \}}\OperatorTok{,}\NormalTok{ \{ }\DataTypeTok{signWith}\OperatorTok{:}\NormalTok{ author \})}\OperatorTok{;}
\end{Highlighting}
\end{Shaded}

Only your key can write it --- enforced by the network, not by
convention. For a friend to read your feed they need two facts: your
\texttt{authorId} and the region you keep it in. Publish those anywhere
(a QR code, an inbox message, a chat), then:

\begin{Shaded}
\begin{Highlighting}[]
\KeywordTok{const}\NormalTok{ theirFeed }\OperatorTok{=}\NormalTok{ \{ }\DataTypeTok{region}\OperatorTok{:} \StringTok{\textquotesingle{}useast\textquotesingle{}}\OperatorTok{,} \DataTypeTok{owner}\OperatorTok{:}\NormalTok{ theirAuthorId}\OperatorTok{,} \DataTypeTok{name}\OperatorTok{:} \StringTok{\textquotesingle{}posts\textquotesingle{}}\NormalTok{ \}}\OperatorTok{;}
\ControlFlowTok{await}\NormalTok{ peer}\OperatorTok{.}\FunctionTok{sub}\NormalTok{(theirFeed}\OperatorTok{,}\NormalTok{ (env) }\KeywordTok{=\textgreater{}} \FunctionTok{renderPost}\NormalTok{(env}\OperatorTok{.}\AttributeTok{message}\NormalTok{)}\OperatorTok{,}\NormalTok{ \{ }\DataTypeTok{since}\OperatorTok{:} \StringTok{\textquotesingle{}all\textquotesingle{}}\NormalTok{ \})}\OperatorTok{;}
\end{Highlighting}
\end{Shaded}

\subsubsection{4.3 An inbox}\label{an-inbox}

\begin{Shaded}
\begin{Highlighting}[]
\KeywordTok{const}\NormalTok{ inboxOf }\OperatorTok{=}\NormalTok{ (authorId) }\KeywordTok{=\textgreater{}}\NormalTok{ (\{ }\DataTypeTok{region}\OperatorTok{:} \StringTok{\textquotesingle{}useast\textquotesingle{}}\OperatorTok{,} \DataTypeTok{owner}\OperatorTok{:}\NormalTok{ authorId}\OperatorTok{,} \DataTypeTok{name}\OperatorTok{:} \StringTok{\textquotesingle{}inbox\textquotesingle{}}\OperatorTok{,} \DataTypeTok{write}\OperatorTok{:} \StringTok{\textquotesingle{}open\textquotesingle{}}\NormalTok{ \})}\OperatorTok{;}

\CommentTok{// anyone sends:}
\ControlFlowTok{await}\NormalTok{ peer}\OperatorTok{.}\FunctionTok{pub}\NormalTok{(}\FunctionTok{inboxOf}\NormalTok{(theirAuthorId)}\OperatorTok{,}\NormalTok{ \{ }\DataTypeTok{from}\OperatorTok{:}\NormalTok{ author}\OperatorTok{.}\AttributeTok{authorId}\OperatorTok{,} \DataTypeTok{text}\OperatorTok{:} \StringTok{\textquotesingle{}hi!\textquotesingle{}}\NormalTok{ \}}\OperatorTok{,}
\NormalTok{               \{ }\DataTypeTok{signWith}\OperatorTok{:}\NormalTok{ author \})}\OperatorTok{;}
\CommentTok{// the owner reads:}
\ControlFlowTok{await}\NormalTok{ peer}\OperatorTok{.}\FunctionTok{sub}\NormalTok{(}\FunctionTok{inboxOf}\NormalTok{(author}\OperatorTok{.}\AttributeTok{authorId}\NormalTok{)}\OperatorTok{,}\NormalTok{ (env) }\KeywordTok{=\textgreater{}} \FunctionTok{showDM}\NormalTok{(env}\OperatorTok{.}\AttributeTok{message}\NormalTok{)}\OperatorTok{,}\NormalTok{ \{ }\DataTypeTok{since}\OperatorTok{:} \StringTok{\textquotesingle{}all\textquotesingle{}}\NormalTok{ \})}\OperatorTok{;}
\end{Highlighting}
\end{Shaded}

The envelope's \texttt{signerPubkey} tells the owner who \emph{really}
sent each message (the \texttt{from} field inside the payload is just a
courtesy copy --- trust the signature, display the field).

\subsubsection{4.4 Who's here? (presence)}\label{whos-here-presence}

\begin{Shaded}
\begin{Highlighting}[]
\KeywordTok{const}\NormalTok{ presence }\OperatorTok{=}\NormalTok{ \{ }\DataTypeTok{region}\OperatorTok{:} \StringTok{\textquotesingle{}useast\textquotesingle{}}\OperatorTok{,} \DataTypeTok{name}\OperatorTok{:} \StringTok{\textquotesingle{}myapp:presence\textquotesingle{}}\NormalTok{ \}}\OperatorTok{;}
\KeywordTok{const}\NormalTok{ here }\OperatorTok{=} \KeywordTok{new} \BuiltInTok{Map}\NormalTok{()}\OperatorTok{;}   \CommentTok{// authorId {-}\textgreater{} last heartbeat}

\PreprocessorTok{setInterval}\NormalTok{(() }\KeywordTok{=\textgreater{}}\NormalTok{ peer}\OperatorTok{.}\FunctionTok{pub}\NormalTok{(presence}\OperatorTok{,}\NormalTok{ \{ }\DataTypeTok{at}\OperatorTok{:} \BuiltInTok{Date}\OperatorTok{.}\FunctionTok{now}\NormalTok{() \}}\OperatorTok{,}\NormalTok{ \{ }\DataTypeTok{signWith}\OperatorTok{:}\NormalTok{ author \})}\OperatorTok{,} \DecValTok{30\_000}\NormalTok{)}\OperatorTok{;}

\ControlFlowTok{await}\NormalTok{ peer}\OperatorTok{.}\FunctionTok{sub}\NormalTok{(presence}\OperatorTok{,}\NormalTok{ (env) }\KeywordTok{=\textgreater{}}\NormalTok{ \{}
  \ControlFlowTok{if}\NormalTok{ (env}\OperatorTok{.}\AttributeTok{signerPubkey}\NormalTok{) here}\OperatorTok{.}\FunctionTok{set}\NormalTok{(env}\OperatorTok{.}\AttributeTok{signerPubkey}\OperatorTok{,}\NormalTok{ env}\OperatorTok{.}\AttributeTok{ts}\NormalTok{)}\OperatorTok{;}
\NormalTok{\})}\OperatorTok{;}
\CommentTok{// online = heartbeat within the last \textasciitilde{}90s}
\KeywordTok{const}\NormalTok{ online }\OperatorTok{=}\NormalTok{ [}\OperatorTok{...}\NormalTok{here]}\OperatorTok{.}\FunctionTok{filter}\NormalTok{(([}\OperatorTok{,}\NormalTok{ t]) }\KeywordTok{=\textgreater{}} \BuiltInTok{Date}\OperatorTok{.}\FunctionTok{now}\NormalTok{() }\OperatorTok{{-}}\NormalTok{ t }\OperatorTok{\textless{}} \DecValTok{90\_000}\NormalTok{)}\OperatorTok{;}
\end{Highlighting}
\end{Shaded}

Presence is just a topic everyone heartbeats into. Messages age out of
the cache on their own; nobody cleans up.

Two details of this recipe are deliberate, and both have been gotten
wrong in a shipped app: the subscription is \textbf{live-only} (no
\texttt{since} --- replaying hours of stale heartbeats is noise), and
recency reads \textbf{\texttt{env.ts}} --- the \emph{publish} time ---
never the local clock at delivery. If you do replay history on a
presence-like topic, a heartbeat that \emph{arrives} now may have been
\emph{published} hours ago; stamping arrival time resurrects everyone
who was online in the last day. Keep the latest \texttt{env.ts} per
author and compare against that.

\subsubsection{4.5 Taking it back (kill)}\label{taking-it-back-kill}

\begin{Shaded}
\begin{Highlighting}[]
\KeywordTok{const}\NormalTok{ msgId }\OperatorTok{=} \ControlFlowTok{await}\NormalTok{ peer}\OperatorTok{.}\FunctionTok{pub}\NormalTok{(}\FunctionTok{room}\NormalTok{(}\StringTok{\textquotesingle{}general\textquotesingle{}}\NormalTok{)}\OperatorTok{,}\NormalTok{ \{ }\DataTypeTok{text}\OperatorTok{:} \StringTok{\textquotesingle{}typo galore\textquotesingle{}}\NormalTok{ \}}\OperatorTok{,}\NormalTok{ \{ }\DataTypeTok{signWith}\OperatorTok{:}\NormalTok{ author \})}\OperatorTok{;}
\ControlFlowTok{await}\NormalTok{ peer}\OperatorTok{.}\FunctionTok{kill}\NormalTok{(}\FunctionTok{room}\NormalTok{(}\StringTok{\textquotesingle{}general\textquotesingle{}}\NormalTok{)}\OperatorTok{,}\NormalTok{ msgId}\OperatorTok{,}\NormalTok{ \{ }\DataTypeTok{signWith}\OperatorTok{:}\NormalTok{ author \})}\OperatorTok{;}   \CommentTok{// same author key!}
\end{Highlighting}
\end{Shaded}

Every subscriber's handler receives
\texttt{\{\ deleted:\ true,\ msgId\ \}} --- drop it from your UI.
Honesty clause: retraction is cooperative, not magical. A device that
already displayed the message has, well, displayed it. And an anonymous
message can never be killed --- no key, no proof it's yours.

\subsubsection{4.6 Sharing files
(std/chunk)}\label{sharing-files-stdchunk}

Single publishes over \textasciitilde15 KB are unreliable on real
browser channels, so don't send big payloads --- chunk them. The kernel
ships the helper:

\begin{Shaded}
\begin{Highlighting}[]
\ImportTok{import}\NormalTok{ \{ publishChunkedBytes}\OperatorTok{,}\NormalTok{ receiveChunkedBytes \} }\ImportTok{from} \StringTok{\textquotesingle{}@axona/protocol/std/chunk.js\textquotesingle{}}\OperatorTok{;}

\KeywordTok{const}\NormalTok{ drop }\OperatorTok{=}\NormalTok{ \{ }\DataTypeTok{region}\OperatorTok{:} \StringTok{\textquotesingle{}useast\textquotesingle{}}\OperatorTok{,} \DataTypeTok{name}\OperatorTok{:} \StringTok{\textquotesingle{}myapp:drops\textquotesingle{}}\NormalTok{ \}}\OperatorTok{;}

\ControlFlowTok{await} \FunctionTok{publishChunkedBytes}\NormalTok{(peer}\OperatorTok{,}\NormalTok{ fileBytes }\CommentTok{/* Uint8Array */}\OperatorTok{,}\NormalTok{ \{}
  \DataTypeTok{topic}\OperatorTok{:}\NormalTok{ drop}\OperatorTok{,} \DataTypeTok{signWith}\OperatorTok{:}\NormalTok{ author}\OperatorTok{,}
  \DataTypeTok{meta}\OperatorTok{:}\NormalTok{ \{ }\DataTypeTok{filename}\OperatorTok{:} \StringTok{\textquotesingle{}demo.png\textquotesingle{}}\OperatorTok{,} \DataTypeTok{mime}\OperatorTok{:} \StringTok{\textquotesingle{}image/png\textquotesingle{}}\NormalTok{ \}}\OperatorTok{,}
\NormalTok{\})}\OperatorTok{;}

\ControlFlowTok{await} \FunctionTok{receiveChunkedBytes}\NormalTok{(peer}\OperatorTok{,}\NormalTok{ drop}\OperatorTok{,}\NormalTok{ \{}
  \DataTypeTok{onComplete}\OperatorTok{:}\NormalTok{ (\{ bytes}\OperatorTok{,}\NormalTok{ meta \}) }\KeywordTok{=\textgreater{}} \FunctionTok{showImage}\NormalTok{(bytes}\OperatorTok{,}\NormalTok{ meta)}\OperatorTok{,}
  \DataTypeTok{onProgress}\OperatorTok{:}\NormalTok{ (\{ received}\OperatorTok{,}\NormalTok{ total \}) }\KeywordTok{=\textgreater{}}\NormalTok{ bar}\OperatorTok{.}\FunctionTok{set}\NormalTok{(received }\OperatorTok{/}\NormalTok{ total)}\OperatorTok{,}
\NormalTok{\})}\OperatorTok{;}
\end{Highlighting}
\end{Shaded}

\subsubsection{4.7 Direct messages (device to
device)}\label{direct-messages-device-to-device}

Pub/sub addresses \emph{audiences}. To address one specific connected
device:

\begin{Shaded}
\begin{Highlighting}[]
\NormalTok{peer}\OperatorTok{.}\FunctionTok{onMessage}\NormalTok{(}\KeywordTok{async}\NormalTok{ (fromNodeId}\OperatorTok{,}\NormalTok{ msg) }\KeywordTok{=\textgreater{}}\NormalTok{ (\{ }\DataTypeTok{pong}\OperatorTok{:}\NormalTok{ msg}\OperatorTok{.}\AttributeTok{ping}\NormalTok{ \}))}\OperatorTok{;}   \CommentTok{// respond}
\KeywordTok{const}\NormalTok{ reply }\OperatorTok{=} \ControlFlowTok{await}\NormalTok{ peer}\OperatorTok{.}\FunctionTok{send}\NormalTok{(someNodeId}\OperatorTok{,}\NormalTok{ \{ }\DataTypeTok{ping}\OperatorTok{:} \DecValTok{1}\NormalTok{ \})}\OperatorTok{;}            \CommentTok{// RPC}
\NormalTok{peer}\OperatorTok{.}\FunctionTok{notify}\NormalTok{(someNodeId}\OperatorTok{,}\NormalTok{ \{ }\DataTypeTok{fyi}\OperatorTok{:} \KeywordTok{true}\NormalTok{ \})}\OperatorTok{;}                            \CommentTok{// fire{-}and{-}forget}
\end{Highlighting}
\end{Shaded}

Node IDs come from \texttt{peer.peers()} / \texttt{peer.onPeerJoin}.
Remember: a node ID is a \emph{session}, not a person --- for
person-to-person, use an inbox (§4.3).

\subsubsection{4.8 How busy is this
topic?}\label{how-busy-is-this-topic}

\begin{Shaded}
\begin{Highlighting}[]
\ImportTok{import}\NormalTok{ \{ metricTopic}\OperatorTok{,}\NormalTok{ deriveTopicId \} }\ImportTok{from} \StringTok{\textquotesingle{}@axona/protocol\textquotesingle{}}\OperatorTok{;}

\KeywordTok{const}\NormalTok{ snap }\OperatorTok{=} \ControlFlowTok{await}\NormalTok{ peer}\OperatorTok{.}\FunctionTok{metrics}\NormalTok{(topic)}\OperatorTok{;}        \CommentTok{// one{-}shot \{ subscribers, current\_count, ... \}}

\ControlFlowTok{await}\NormalTok{ peer}\OperatorTok{.}\FunctionTok{sub}\NormalTok{(}\FunctionTok{metricTopic}\NormalTok{(}\ControlFlowTok{await} \FunctionTok{deriveTopicId}\NormalTok{(topic))}\OperatorTok{,}\NormalTok{ (env) }\KeywordTok{=\textgreater{}}\NormalTok{ \{}
  \KeywordTok{const}\NormalTok{ m }\OperatorTok{=} \BuiltInTok{JSON}\OperatorTok{.}\FunctionTok{parse}\NormalTok{(env}\OperatorTok{.}\AttributeTok{message}\NormalTok{)}\OperatorTok{;}
\NormalTok{  badge}\OperatorTok{.}\AttributeTok{textContent} \OperatorTok{=} \VerbatimStringTok{\textasciigrave{}}\SpecialCharTok{$\{}\NormalTok{m}\OperatorTok{.}\AttributeTok{subscribers}\SpecialCharTok{\}}\VerbatimStringTok{ watching\textasciigrave{}}\OperatorTok{;}
\NormalTok{\}}\OperatorTok{,}\NormalTok{ \{ }\DataTypeTok{since}\OperatorTok{:} \StringTok{\textquotesingle{}latest\textquotesingle{}}\NormalTok{ \})}\OperatorTok{;}                       \CommentTok{// live{-}updating counter}
\end{Highlighting}
\end{Shaded}

Metrics are \textbf{published on demand}: your subscription is what
switches them on, and the root answers \textbf{immediately} --- the
first snapshot arrives at routing latency (\textasciitilde0.3 s on
testnet), with or before your data replay, whether or not anyone has
ever published (then one every \textasciitilde20 s while you stay
subscribed). Two habits worth keeping:

\begin{itemize}
\tightlist
\item
  \textbf{Treat silence as \emph{unknown}, not zero.} Until the first
  snapshot arrives, ``no activity'' isn't an answer yet --- a
  \texttt{current\_count:\ 0} snapshot is.
\item
  \textbf{Keep the subscription open.} A sub torn down seconds after it
  starts can still miss the answer under churn; the stream is the
  primitive, not a one-shot read.
\end{itemize}

Counts are advisory --- decorate with them, never authorize with them.

\subsubsection{4.9 Remembering your users}\label{remembering-your-users}

The complete persistence story for most apps is one line you've already
written:
\texttt{createAuthorIdentity(\{\ persistAs:\ \textquotesingle{}myapp:author\textquotesingle{}\ \})}.
Node identity? Re-mint each run. Subscriptions? Re-subscribe on startup
(with \texttt{since}, you'll catch up on what you missed). Messages? The
network holds the recent window; anything your app must keep beyond
\textasciitilde24 h, store like any other app data.

\subsubsection{4.10 Playing nicely with other
apps}\label{playing-nicely-with-other-apps}

Topics are shared space --- another app can subscribe to your topic. If
you want your messages legible to it (and its to you), use the standard
body:

\begin{Shaded}
\begin{Highlighting}[]
\ImportTok{import}\NormalTok{ \{ makeMessage}\OperatorTok{,}\NormalTok{ readMessage \} }\ImportTok{from} \StringTok{\textquotesingle{}@axona/protocol/std/message.js\textquotesingle{}}\OperatorTok{;}

\ControlFlowTok{await}\NormalTok{ peer}\OperatorTok{.}\FunctionTok{pub}\NormalTok{(topic}\OperatorTok{,} \FunctionTok{makeMessage}\NormalTok{(}\StringTok{\textquotesingle{}hello\textquotesingle{}}\OperatorTok{,}\NormalTok{ \{ }\DataTypeTok{mood}\OperatorTok{:} \StringTok{\textquotesingle{}sunny\textquotesingle{}}\NormalTok{ \})}\OperatorTok{,}\NormalTok{ \{ }\DataTypeTok{signWith}\OperatorTok{:}\NormalTok{ author \})}\OperatorTok{;}
\KeywordTok{const}\NormalTok{ \{ text \} }\OperatorTok{=} \FunctionTok{readMessage}\NormalTok{(env}\OperatorTok{.}\AttributeTok{message}\NormalTok{)}\OperatorTok{;}    \CommentTok{// tolerant of other apps\textquotesingle{} shapes}
\end{Highlighting}
\end{Shaded}

\subsubsection{4.11 If your author is a bot, say
so}\label{if-your-author-is-a-bot-say-so}

\begin{Shaded}
\begin{Highlighting}[]
\ControlFlowTok{await}\NormalTok{ peer}\OperatorTok{.}\FunctionTok{setAuthorClass}\NormalTok{(}\StringTok{\textquotesingle{}agent\textquotesingle{}}\OperatorTok{,}\NormalTok{ \{ }\DataTypeTok{signWith}\OperatorTok{:}\NormalTok{ author}\OperatorTok{,} \DataTypeTok{label}\OperatorTok{:} \StringTok{\textquotesingle{}ticker{-}bot\textquotesingle{}}\NormalTok{ \})}\OperatorTok{;}
\end{Highlighting}
\end{Shaded}

A voluntary, signed declaration that this key is operated by software.
Readers resolve it with \texttt{getAuthorClass(authorId)} and can badge,
rank, or filter. It gates nothing --- it's honesty infrastructure, and
well-behaved bots (and AI-built apps!) should use it.

\subsection{5. Things that will bite
you}\label{things-that-will-bite-you}

Learn these here rather than in production. The first five are the
classic protocol traps; the last five have each already bitten a real,
shipped Axona application.

\textbf{1. Two descriptors, one character apart, are two different
topics.} No error, no warning --- just silence. The classic: publisher
says
\texttt{\{\ region:\ \textquotesingle{}useast\textquotesingle{},\ name:\ \textquotesingle{}lobby\textquotesingle{}\ \}},
subscriber says
\texttt{\{\ name:\ \textquotesingle{}lobby\textquotesingle{}\ \}} from a
laptop in Berlin, and its defaulted region is \texttt{eucentral}. One
shared \texttt{topics.js} with functions like \texttt{room(name)} ends
this bug forever.

\textbf{2. There is no ack --- stop waiting for one.} \texttt{pub()}
resolves when the message is \emph{sent}, not delivered, and no receipt
ever comes (it would stitch your connection to your authorship --- an
identity leak Axona refuses on principle). The idiom: you're subscribed
to what you publish; when your own \texttt{msgId} comes back through
your handler, it went through. The kernel meanwhile retries quietly on
your behalf.

\textbf{3. Publishing before \texttt{ready()} is a coin flip.} The mesh
takes a few seconds to warm after \texttt{start()}.
\texttt{await\ peer.ready()} --- that's what it's for. (Your first
publish would \emph{probably} survive anyway; the kernel bursts a
newcomer's first messages. ``Probably'' is not a UX.)

\textbf{4. Messages are not rows.} The network keeps a rolling
\textasciitilde24-hour window per topic, and open topics quota each
author's share of it. Build for a living stream: late joiners see recent
history, not the archive. Need an archive? That's your app's job (and
re-publishing refreshes a message you want kept warm).

\textbf{5. 15 KB.} Above it, real-world WebRTC channels start dropping
frames even though the hard cap is 256 KB. \texttt{std/chunk} exists so
you never think about this again.

\textbf{6. Delivery time is not publish time.} A subscription with
\texttt{since} replays history: a message that arrives \emph{now} may
have been published hours ago. Any recency logic --- ``last seen'',
online lists, freshness sorting --- reads \texttt{env.ts} (the publish
time), never \texttt{Date.now()} at delivery. The observed failure: an
app stamped arrival time on replayed heartbeats, and every user from the
past day appeared ``online'' for 90 seconds after each launch.

\textbf{7. Switching users is not reconnecting.} One peer serves any
number of authors --- authorship is per-call
(\texttt{\{\ signWith\ \}}), and swapping the active persona should
touch \emph{nothing} about the connection. The observed failure: an app
tore down and rebuilt its peer on every persona switch, churning the
mesh and killing live connections mid-negotiation for zero benefit.

\textbf{8. Sharing a topic means sharing the whole descriptor.} An
invite, a directory advertisement, a QR code, a pasted link --- whatever
carries a topic to another user must carry
\texttt{\{\ region,\ owner,\ name,\ write\ \}} \emph{complete}. The
observed failure: a discovery ticker advertised owned topics by name and
region only; joiners derived a different (ownerless) topic ID and landed
in a working-looking, permanently empty room.

\textbf{9. When derivation fails, fail.} If the kernel rejects a
descriptor, show the error and skip the topic. The observed failure: an
app caught the throw and substituted a locally invented ID ``so the UI
kept working'' --- which converted a visible error into invisible,
permanent non-delivery that no log would ever explain.

\textbf{10. Local dev: bind the dev server to IPv4.} Vite's default
\texttt{localhost} binding lands on IPv6 \texttt{::1}, and Firefox
gathers \textbf{zero} ICE candidates on a page served from a
\texttt{::1} origin --- every mesh dial fails with a misleading ``your
TURN server appears to be broken'' console error, while Chromium works.
One line ends it:
\texttt{server:\ \{\ host:\ \textquotesingle{}127.0.0.1\textquotesingle{}\ \}}.
(Bundler note: alias \texttt{node-datachannel} and
\texttt{node-datachannel/polyfill} to an empty stub for browser builds,
and clear the bundler's dependency cache after changing the kernel pin.)

\subsection{6. Errors, when they happen}\label{errors-when-they-happen}

Everything the kernel throws is an \texttt{AxonaError} with a stable
\texttt{.code} --- switch on codes, not message text:

{\def\LTcaptype{none} % do not increment counter
\begin{longtable}[]{@{}
  >{\raggedright\arraybackslash}p{(\linewidth - 4\tabcolsep) * \real{0.3333}}
  >{\raggedright\arraybackslash}p{(\linewidth - 4\tabcolsep) * \real{0.3333}}
  >{\raggedright\arraybackslash}p{(\linewidth - 4\tabcolsep) * \real{0.3333}}@{}}
\toprule\noalign{}
\begin{minipage}[b]{\linewidth}\raggedright
Code
\end{minipage} & \begin{minipage}[b]{\linewidth}\raggedright
You did
\end{minipage} & \begin{minipage}[b]{\linewidth}\raggedright
Do instead
\end{minipage} \\
\midrule\noalign{}
\endhead
\bottomrule\noalign{}
\endlastfoot
\texttt{PUBLISH\_NO\_PUBLISH\_IDENTITY} & \texttt{pub} without
\texttt{signWith} & pass an author, or \texttt{ANONYMOUS} explicitly \\
\texttt{WRITE\_POLICY\_VIOLATION} & wrote to someone else's owned topic
& only the owner key writes; use their inbox \\
\texttt{TOPIC\_REGION\_REQUIRED} & derived an open topic with no region
anywhere & name a \texttt{region} \\
\texttt{PUBLISH\_PAYLOAD\_TOO\_LARGE} & \textgreater256 KB publish &
\texttt{std/chunk} \\
\texttt{KILL\_SIGN\_FAILED} & killed with a different key & retract with
the key that signed \\
WS close \texttt{4426} & version mismatch with the bridge & reinstall
the pinned kernel tag \\
\end{longtable}
}

Not errors, by design: \texttt{pull()} returning \texttt{null} (gone or
never existed) and \texttt{kill()} resolving \texttt{\{\ ok:\ false\ \}}
(nothing to retract).

\subsection{7. Limits worth knowing}\label{limits-worth-knowing}

{\def\LTcaptype{none} % do not increment counter
\begin{longtable}[]{@{}
  >{\raggedright\arraybackslash}p{(\linewidth - 2\tabcolsep) * \real{0.5000}}
  >{\raggedright\arraybackslash}p{(\linewidth - 2\tabcolsep) * \real{0.5000}}@{}}
\toprule\noalign{}
\begin{minipage}[b]{\linewidth}\raggedright
Thing
\end{minipage} & \begin{minipage}[b]{\linewidth}\raggedright
Value
\end{minipage} \\
\midrule\noalign{}
\endhead
\bottomrule\noalign{}
\endlastfoot
Reliable message size & \textbf{15 KB} (chunk above; 256 KB absolute) \\
Message lifetime & \textasciitilde24 h (48 h max; re-publish to
refresh) \\
History for late joiners & the cached queue --- recent, bounded, not
forever \\
One author's share of an open topic & quota-bounded (no flooding) \\
Clock skew tolerance on live publishes & ±5 minutes (keep device clocks
sane) \\
\end{longtable}
}

\subsection{8. Going live}\label{going-live}

\begin{itemize}
\tightlist
\item
  \textbf{Point at the right network.} Both live networks run this
  guide's 4.x line: \textbf{testnet} (\texttt{wss://testnet.axona.net})
  tracks the newest kernel, \textbf{production}
  (\texttt{wss://bridge.axona.net}, with a federated west sibling) runs
  the most recently promoted release, typically one behind. Develop
  against testnet; deploy against production. A genuinely mismatched
  kernel is refused loudly (close 4426) --- match your kernel pin to
  your bridge.
\item
  \textbf{HTTPS is mandatory} in the browser (\texttt{crypto.subtle}).
\item
  \textbf{Prove it with two clients before you call it done.} A
  distributed app can pass every single-client test while being
  completely wrong --- your own client shows you your own state, not the
  network's. The minimum gate: two clients in one topic exchanging
  messages; a kill on one disappearing on the other; a topic shared
  through your app's own invite/ad/QR surface opening with
  \emph{content} (not a clean-looking empty room --- that's a partial
  descriptor); a fresh client's presence view showing only the currently
  active; and the whole gate run in both a Chromium browser and Firefox.
  Every field failure in §5's entries 6--10 was invisible to one client.
\item
  \textbf{Surface your versions.} Show the kernel version somewhere
  findable (\texttt{KERNEL\_VERSION} import) --- future-you debugging a
  mixed fleet will be grateful.
\item
  \textbf{Durability beyond browser tabs.} An app whose topics matter
  when no user is online wants a small always-on peer that
  \texttt{host()}s them --- that's a \emph{relay}, a five-minute Node
  deployment. When you reach that point, the
  \href{Axona-Services-Guide-v4.27.1.md}{Services Guide} is yours; not
  before.
\end{itemize}

\subsection{\texorpdfstring{9. Under the hood ---
\emph{optional}}{9. Under the hood --- optional}}\label{under-the-hood-optional}

You can build everything above without this section. For the curious:

\textbf{The mesh.} Peers connect via a bridge (a meeting point, not a
server --- it relays introductions and can pass messages, but owns
nothing), then directly to each other over WebRTC. Each peer maintains a
small routing table; the network is navigable without anyone holding a
map. It self-heals: dead connections are detected by heartbeat and
routed around, no configuration, no ops.

\textbf{The tree.} Every topic ID is an address in the same space as
peer IDs. The live peer \emph{closest} to a topic's address
automatically becomes its \textbf{root} --- the coordinator that stamps,
caches, and fans out that topic's messages. Nobody elects it; proximity
appoints it. Popular topics grow a distribution tree (the root delegates
batches of subscribers to child relays), so no node ever serves more
than a bounded number directly.

\textbf{Durability.} The root continuously copies its cache to the peers
next in line; when a root vanishes (a laptop closes), a warm backup
notices within a minute and takes over, and every subscriber's periodic
renewal re-attaches automatically. A peer that leaves \emph{gracefully}
hands every topic it was coordinating to an heir before it goes. And
whenever a topic's coordinator changes hands, the old and new holders
reconcile to the \textbf{union} of their history --- so a newcomer
asking for ``everything'' gets the whole timeline, not whichever half
its node happened to hold. Your app sees, at worst, a brief pause. Your
messages are also retried by your own peer until it observes them
delivered. All of this is why the guide keeps saying \emph{you don't
manage the network} --- it manages itself, aggressively.

\textbf{Regions.} The first byte of every ID is one of 192 real
geographic cells, so a region's topics naturally live on that region's
peers --- locality without a CDN. (A stricter mode, where a topic may
\emph{only} be served in its own region, is built and waiting for
network density to justify switching on.)

The full story --- with diagrams --- is the
\href{../architecture/Axona-Architecture.tex}{Axona Architecture} note.

\subsection{10. Appendix}\label{appendix}

\subsubsection{Migrating from the v2/v3
lines}\label{migrating-from-the-v2v3-lines}

Only relevant if you have code from early 2026. The v3 flag-day replaced
\texttt{deriveIdentity} with the two factories, string topics with
descriptors, and \texttt{publisher}/\texttt{sign:false} with per-publish
\texttt{signWith}; v4 changed nothing in the app API --- it rebuilt
reliability underneath (single-root axon trees, cohort replication,
cold-publish burst, nearest-replica reads). \texttt{unpub} was removed
and \texttt{touch} retired (re-publish to refresh a message instead). If
a symbol you remember is missing from the
\href{Axona-API-Reference-v4.27.1.md}{API Reference}, it didn't survive
v3 --- the reference is the source of truth.

\begin{center}\rule{0.5\linewidth}{0.5pt}\end{center}

\emph{Axona Programmer Guide · kernel 4.27.1 · testnet · 2026-07-17}

\end{document}
